\hypertarget{discussion}{%
\section{4. Discussion}\label{discussion}}

\hypertarget{interpretation-of-the-main-findings}{%
\subsection{4.1. Interpretation of the Main
Findings}\label{interpretation-of-the-main-findings}}

The successful application of Machine Learning (ML) to the monitoring of
Equatorial Plasma Bubbles (EPBs) signifies a paradigm shift from
labor-intensive, subjective manual observation to automated, highly
scalable analytical frameworks. The findings demonstrate that deep
learning and advanced feature extraction algorithms do not merely
replicate human classification but can identify subtle spatial
gradients, textures, and intensity variations that are often visually
ambiguous to the human eye. Furthermore, the automated extraction of
morphological parameters---such as the geometric center of the western
wall and the northernmost point---has proven robust enough to generate
the first solar-cycle-scale statistical profiles of EPBs, revealing
complex variations tied to solar and geomagnetic activity that were
previously impossible to compile efficiently.

As evidenced by the compiled results, there is a dominant methodological
reliance on optical imagery coupled with DCNN architectures. However,
the overarching limitations reported across these studies consistently
point toward a strong, inherent susceptibility to environmental
contamination (e.g., cloud cover) and the systemic challenge of handling
heavily imbalanced datasets driven by natural solar cycle variations.
These structural trends establish clear performance trade-offs and
operational boundary conditions, which are synthesized in detail in the
subsequent sections.

\hypertarget{comparison-between-machine-learning-techniques}{%
\subsection{4.2. Comparison Between Machine Learning
Techniques}\label{comparison-between-machine-learning-techniques}}

A critical analysis of the deployed architectures reveals a distinct
trade-off between computational complexity and operational efficiency.
Deep Convolutional Neural Networks (DCNNs) such as ResNet, Mask R-CNN,
and YOLO are highly effective at capturing the non-rigid, branching
morphologies of EPBs and performing pixel-level semantic segmentation.
However, increasing the depth of these networks often yields only
marginal improvements in accuracy while significantly increasing
processing latency, making them computationally expensive.

Conversely, hybrid models integrating dimensionality reduction with
traditional classifiers---such as Two-Dimensional Principal Component
Analysis (2DPCA) or Singular Value Decomposition (SVD) combined with
Support Vector Machines (SVM) or Random Forests (RF)---offer a highly
efficient alternative. By directly processing the 2D image matrix and
eliminating redundant features, these techniques dramatically reduce
trained model size (by up to 70\%) and training time without
compromising detection accuracy. Additionally, the use of ensemble
techniques, such as bootstrapping CNNs with probability averaging or
mode voting, mitigates the overconfidence of individual sub-models and
enhances robustness against ambiguous cases. Finally, the integration of
Explainable AI (XAI) frameworks, specifically SHAP and LIME, resolves
the ``black-box'' dilemma of deep learning by proving that these models
base their classifications on genuine physical features---like the
vertical extent and width of dark plumes---rather than on background
artifacts.

\hypertarget{impact-of-different-data-types}{%
\subsection{4.3. Impact of Different Data
Types}\label{impact-of-different-data-types}}

The efficacy of ML models is intrinsically linked to the data modality
utilized. Optical All-Sky Imagers (ASI) remain the primary data source
due to their wide field of view and operational ability to capture
fine-scale 2D morphologies and bifurcation structures. In contrast,
Global Navigation Satellite System (GNSS) data, utilized through Total
Electron Content (TEC) maps and ROTI indices, offers
weather-independent, continuous global coverage. While detrended TEC
maps successfully identify wave-like perturbations like MSTIDs, their
spatial resolution is generally inferior to ASI for observing the highly
localized, fine structures of EPBs. Very High Frequency (VHF) radar
provides exceptionally high spatial and temporal resolution for
observing bottom-side F-layer boundary irregularities, but it is
constrained by a localized geographical footprint. The complementary
nature of these data types suggests that relying on a single instrument
limits the scope of monitoring systems.

\hypertarget{recurring-limitations-in-the-studies}{%
\subsection{4.4. Recurring Limitations in the
Studies}\label{recurring-limitations-in-the-studies}}

Despite high reported accuracies, several recurring limitations persist
across the literature. Foremost is the \textbf{severe susceptibility to
environmental contamination}. While ASI provides superior resolution,
its inherent reliance on optical visibility means that tropospheric
weather conditions---such as heavy cloud cover, lunar illumination, or
artificial light pollution---can obscure EPB signatures. This
drastically reduces the reliability of optical-based models and
frequently leads to misclassifications.

Another critical limitation is the \textbf{imbalance in training
datasets}. EPB occurrence is highly correlated with solar activity,
resulting in datasets skewed heavily toward high solar activity years.
Consequently, models tend to become biased toward the majority class,
performing optimally during solar maxima or geomagnetically quiet
periods, but struggling during solar minima or severe geomagnetic storms
where sample sizes are much smaller.

Furthermore, because these are predominantly supervised learning models,
they rely entirely on \textbf{manual ground-truth annotations}. The
subjective nature of manually distinguishing between fluid fluid
transitions (e.g., calm vs.~dynamic episodes) or faint EPBs from clouds
establishes an artificial ceiling on model accuracy, as human error is
inadvertently embedded into the training data.

\hypertarget{challenges-in-the-field}{%
\subsection{4.5. Challenges in the
Field}\label{challenges-in-the-field}}

A major challenge lies in the \textbf{generalizability of the models}.
Models trained on data from a specific geographical sector (e.g.,
Brazil, China, or Thailand) or a specific imager configuration may
experience degraded performance when applied to different locations due
to varying camera characteristics, viewing geometries, and regional EPB
morphologies. Additionally, balancing real-time operational processing
speeds with the computational demands of deep learning algorithms
remains a hurdle for deploying these systems on low-cost, edge-computing
devices.

\hypertarget{gaps-in-the-literature}{%
\subsection{4.6. Gaps in the Literature}\label{gaps-in-the-literature}}

Current research is heavily skewed toward reactive classification and
detection. There is a \textbf{stark lack of predictive models} capable
of forecasting the initiation, growth rate, and future trajectory of
EPBs before they fully manifest. Furthermore, while researchers often
use multiple instruments (e.g., ASI and GNSS) to validate findings
independently, there is an \textbf{absence of true multimodal deep
learning frameworks} that ingest and fuse optical, radar, and TEC data
simultaneously within a single neural network architecture to overcome
individual instrumental blind spots.

\hypertarget{opportunities-for-future-work}{%
\subsection{4.7. Opportunities for Future
Work}\label{opportunities-for-future-work}}

Future research should prioritize the development of
\textbf{multi-instrument data fusion models}, integrating ASI, GNSS-TEC,
and satellite in-situ measurements to create robust systems that
maintain high accuracy regardless of tropospheric weather conditions.
Additionally, there is a critical need to transition from detection to
\textbf{spatiotemporal predictive modeling}. By applying sequence-based
networks (e.g., LSTMs or Transformers) to time-series airglow imagery
and geophysical indices, models could proactively forecast EPB zonal
drift and morphological evolution. Finally, XAI should be leveraged not
just for model validation, but for \textbf{scientific discovery}, using
interpretability tools to uncover previously unnoticed physical
relationships and environmental drivers governing ionospheric dynamics.
